\LectureAnswerSection

\begin{drillsolution}{t2-11}\textbf{用相同温度确定每个分子的平均平动动能}\par\nobreak
同温约束的是每个分子的平均平动动能。氦和氮气均满足 $\overline\varepsilon=3kT/2$，所以能量之比为 $\boxed{1:1}$。氮气还具有转动内能，但题问平均平动动能，不能把转动自由度加进来。


\par\medskip\Needspace{4\baselineskip}\noindent\textbf{用反平方根关系比较方均根速率}\par\nobreak
由 $\tfrac12m\langle v^2\rangle=3kT/2$ 得方均根速率与分子质量的平方根成反比；分子质量比与摩尔质量比相同。图\ref{fig:t2-11s}在固定温度下画出这一反平方根关系，两个点分别代入氦与氮气的摩尔质量；纵轴比较的是方均根速率，并非每个分子的实际速率。\fig{fig:t2-11s}{同温时方均根速率按摩尔质量平方根的倒数变化}{\begin{tikzpicture}[x=1.1cm,y=1cm]
\draw[->](0,0)--(5.3,0)node[right]{$\mu/(7\,\mathrm{g\,mol^{-1}})$};\draw[->](0,0)--(0,3.6)node[above]{$v_{\rm rms}/v_{{\rm rms},N_2}$};
\draw[blue,thick,domain=.4:4.8,samples=100]plot(\x,{2/sqrt(\x)});
\fill(.571429,2.64575)circle(2pt)node[right]{He：$\mu=4$};\fill(4,1)circle(2pt)node[above right]{$N_2$：$\mu=28$};
\draw[dashed](.571429,0)--(.571429,2.64575)--(0,2.64575) (4,0)--(4,1)--(0,1);
\node[left]at(0,1){$1$};\node[left]at(0,2.64575){$\sqrt7$};\end{tikzpicture}}

\[\boxed{\frac{v_{{\mathrm{rms}},\rm He}}{v_{{\mathrm{rms}},\rm N_2}}
=\sqrt{\frac{28}{4}}=\sqrt7}.\]


\par\medskip\Needspace{4\baselineskip}\noindent\textbf{把质量与均方速率的补偿代回压强式，比较数密度}\par\nobreak
再用数密度形式的状态方程 $p=n_NkT$。由图\ref{fig:t2-11s}的速率比代回微观压强式复核：氮分子质量是氦原子的7倍，同温时均方速率却只有氦的 $1/7$，二者在 $m\langle v^2\rangle$ 中恰好抵消。剩下的差别仅是数密度，故
\[\boxed{p_{\mathrm{N_2}}:p_{\mathrm{He}}=2:1}.\]

\end{drillsolution}
\begin{drillsolution}{t2-12}\textbf{只用初末状态的压强体积乘积求共同内能差}\par\nobreak
先从两个端点求与路径无关的内能差。对1mol单原子气体，$U=3pV/2$；$A$ 点 $pV=p_0V_0$，$C$ 点 $pV=4p_0V_0$，所以两条路径都有
\[\boxed{\Delta U=\frac32(4-1)p_0V_0=\frac92p_0V_0}.\]


\par\medskip\Needspace{4\baselineskip}\noindent\textbf{分别用矩形与梯形面积求两条路径的功}\par\nobreak
路径一的 $AB$ 等容，功为零；$BC$ 压强恒为 $2p_0$，故 $W_1=2p_0(2V_0-V_0)=2p_0V_0$。路径二是直线，其功为梯形面积，
\[W_2=\frac{p_0+2p_0}{2}(2V_0-V_0)=\frac32p_0V_0.\]
图\ref{fig:t2-12s}的阴影正是两条路径所围的三角形，可直接验证 $W_1-W_2=p_0V_0/2$。\fig{fig:t2-12s}{两条实际路径共享$A,C$；蓝色三角形只表示两功之差，不是任一路径的全部做功}{\begin{tikzpicture}

\draw[axis](0,0)--(5.2,0)node[right]{$V/V_0$};\draw[axis](0,0)--(0,4.8)node[above]{$p/p_0$};\fill[blue!12](2,2)--(2,4)--(4,4)--cycle;\draw[flow,blue](2,2)--(2,4)--(4,4);\draw[flow,red](2,2)--(4,4);\fill(2,2)circle(.05)node[left]{$A$};\fill(2,4)circle(.05)node[above left]{$B$};\fill(4,4)circle(.05)node[above right]{$C$};\node[below]at(2,0){$1$};\node[below]at(4,0){$2$};\node[left]at(0,2){$1$};\node[left]at(0,4){$2$};\node at(3,-1){$W_1-W_2=p_0V_0/2$};

\end{tikzpicture}}



\par\medskip\Needspace{4\baselineskip}\noindent\textbf{对每条路径相加求热量，并用阴影面积核对热量差}\par\nobreak
分别代入 $Q=\Delta U+W$，得
\[\boxed{Q_1=\frac{13}{2}p_0V_0,\qquad Q_2=6p_0V_0}.\]
路径一在更多体积变化范围内保持较高压强，因此做功更多；末态内能相同，这部分额外输出只能由额外吸热补偿。故 $Q_1-Q_2=W_1-W_2=p_0V_0/2$，与阴影面积完全一致。
\end{drillsolution}
\begin{drillsolution}{t2-13}\textbf{用刚性与绝热条件分别排除体积功和壁面传热}\par\nobreak
取气体与电阻为系统，外部电源不在系统内。图\ref{fig:t2-13q}所示容器刚性，外边界不移动，故 $W_{\text{边界}}=0$；外壁绝热，故 $Q_{\text{壁}}=0$。电源仍可以通过导线输入600J电功，这与绝热条件并不冲突。


\par\medskip\Needspace{4\baselineskip}\noindent\textbf{把导线电功计入边界内的内能增量}\par\nobreak
忽略电阻最终储能变化，输入能量全部成为气体内能。图\ref{fig:t2-13s}中的边界条件与电功项合并为
\[nC_V\Delta T=600\,\mathrm J.\]\fig{fig:t2-13s}{电源通过两根导线向边界内电阻输入能量，绝热壁阻止的是热传递，刚性壁无体积功}{\begin{tikzpicture}

\draw[world,very thick](0,0)rectangle(4.5,3);\draw[guide,rounded corners](-.2,-.2)rectangle(4.7,3.2);\draw[world](1.5,.8)rectangle(2.5,1.2);\draw[world](1.5,1)--(.8,1)--(.8,2.2)--(6,2.2)--(6,1.8);\draw[world](2.5,1)--(6,1)--(6,1.4);\draw[world](5.6,1.8)--(6.4,1.8) (5.8,1.4)--(6.2,1.4);\node[right]at(6.5,1.6){外电源};\draw[flow,red](5.5,2.6)--(3.5,2.6)node[midway,above]{$W_e=600\,\mathrm J$};\node at(2.2,.4){气体与电阻};\node at(2.2,3.8){$Q_{\text{壁}}=0,\quad W_{\text{边界}}=0$};

\end{tikzpicture}}



\par\medskip\Needspace{4\baselineskip}\noindent\textbf{用气体热容求温升，再用固定体积求压强比}\par\nobreak
单原子气体 $C_V=3R/2$、$n=2\,\mathrm{mol}$，因此
\[\boxed{\Delta T=\frac{600}{2(3R/2)}=24.07\,\mathrm K}.\]
体积和气体量都不变，压强比只需按绝对温度比计算：
\[\boxed{\frac{p_f}{p_i}=\frac{300+24.07}{300}=1.0802,\qquad W_{\text{边界}}=0}.\]


\par\medskip\Needspace{4\baselineskip}\noindent\textbf{改变边界时重新命名能量传递，避免重复计入}\par\nobreak
若只把气体作为系统，可以将电阻向气体传递的600J记为热量，最终温升相同；但不能在同一边界内既加600J电功又加600J热量，否则同一能量被计算了两次。
\end{drillsolution}
\begin{drillsolution}{t2-14}\textbf{取两个容器的整体，先确定外部热量与功均为零}\par\nobreak
打开阀门后，整个外边界仍刚性、绝热，故没有外部体积功和热量交换。氦与氮气可以互相传能并扩散，但两者总内能不变。应各用自己的热容求初态内能，而不能先平均温度。


\par\medskip\Needspace{4\baselineskip}\noindent\textbf{分别使用两种气体的热容建立初末内能守恒}\par\nobreak
氦的 $C_{V,\rm He}=3R/2$，氮气按刚性双原子模型取 $C_{V,\rm N_2}=5R/2$。两者各1mol，末态同温 $T_f$，所以
\[\frac32R\times300+\frac52R\times600
=\left(\frac32+\frac52\right)RT_f.\]
图\ref{fig:t2-14s}用整体边界说明：热量在两种组分之间重新分配，外部没有能量流入，且末态每种组分都占据全部空间。\fig{fig:t2-14s}{绝热外边界内由冷热气体交换能量并混合}{\begin{tikzpicture}[scale=.85]
\draw[very thick](0,0)rectangle(7,2.2);\draw[dashed](3.5,0)--(3.5,2.2);\draw[->,red,thick](5,1.2)--(2,1.2)node[midway,above]{能量交换};
\node at(1.7,.5){He，300K，$V$};\node at(5.3,.5){$N_2$，600K，$V$};
\draw[->](3.5,-.3)--(3.5,-1);\draw[very thick](0,-3.6)rectangle(7,-1.4);
\node at(3.5,-2.1){混合后：$T_f=487.5\,\mathrm K$，总体积 $2V$};
\foreach \x/\y in {1/-3,2.4/-2.8,4/-3.1,5.6/-2.9}{\fill[blue](\x,\y)circle(2pt);\draw[red](\x+.4,\y)circle(2pt);}
\node at(3.5,2.8){外边界绝热、刚性：$U_{\text{初}}=U_{\text{末}}$};\end{tikzpicture}}



\par\medskip\Needspace{4\baselineskip}\noindent\textbf{按热容加权求末温并核对冷热之间的范围}\par\nobreak
\[\boxed{T_f=\frac{(3/2)300+(5/2)600}{4}=487.5\,\mathrm K}.\]
末温处于300K与600K之间，且更靠近热容较大的氮气初温。

\par\medskip\Needspace{4\baselineskip}\noindent\textbf{让每种组分占据全部最终空间，再相加分压}\par\nobreak
整个末态体积为 $2V$，总物质的量2mol，因此
\[\boxed{p_f=\frac{2RT_f}{2V}=\frac{487.5R}{V}}.\]
两种组分都分布在整个 $2V$ 空间内，各有分压 $RT_f/(2V)$，相加才是总压。不能把每种气体仍按原来只占 $V$ 计算，否则会把末压重复增大一倍。
\end{drillsolution}
